\documentclass{article}
\usepackage[UTF8]{ctex}
\usepackage{amsmath,amssymb,amsfonts}
\usepackage{enumitem}
\usepackage{geometry}
\geometry{a4paper, margin=1in}
% === 错题本排版优化（不参与题目标识） ===
\setlength{\parindent}{0pt}
\setlength{\parskip}{0.35em plus 0.10em minus 0.08em}
\setlength{\emergencystretch}{3em}
\setlist{itemsep=0.25em, topsep=0.40em, parsep=0pt}
\allowdisplaybreaks[2]
% === 错题本排版优化结束 ===
\begin{document}
\section*{中国人民大学2026年高等代数真题}
\section*{18．中国人民大学2026年高等代数试卷}

\subsection{题目1}
1．（15 分）设 \(A\) 为 \(n\) 阶实矩阵，\(r(A)=r<n\) ，证明：\(A\) 可写成 \(n-r\) 个秩为 \(n-1\) 的 \(n\) 阶矩阵的乘积．

\subsection{题目2}
2．（15 分）设 \(X=\left(\begin{array}{ll}A & B \\ C & D\end{array}\right)\) ，其中 \(A, B, C, D\) 均为 \(n\) 阶实矩阵，且 \(A\) 是可逆对称矩阵，\(B^{\mathrm{T}}=C\) 。

\begin{itemize}
  \item[（1）] （10 分）证明：存在可逆矩阵 \(P\) ，使得 \(P^{\mathrm{T}} X P\) 为分块对角矩阵。
  \item[（2）] （10分）给出 \(\left(\begin{array}{cc}A & B \\ C & D\end{array}\right)\) 正定的充分必要条件．
\end{itemize}

\subsection{题目3}
3．（15 分）已知实数 \(a_{i} \neq 0(i=1,2, \cdots, n)\) ，计算 \(n\) 阶行列式

\[
D_{n}=\left|\begin{array}{cccc}
0 & a_{1}+a_{2} & \cdots & a_{1}+a_{n} \\
a_{2}+a_{1} & 0 & \cdots & a_{2}+a_{n} \\
\vdots & \vdots & & \vdots \\
a_{n}+a_{1} & a_{n}+a_{2} & \cdots & 0
\end{array}\right|
\]

\subsection{题目4}
4．（20 分）设 \(\alpha_{1}=(1,-1,2,1,0), \alpha_{2}=(2,-2,4,-2,0), \alpha_{3}=(3,0,6,-1,1), \alpha_{4}=(0,3,0,0,1)\) ．

\begin{itemize}
  \item[（1）] （10 分）求该向量组的一个极大线性无关组，并将向量组中其他向量用该极大线性无关组线性表示．
  \item[（2）] （10 分）求一个以该极大线性无关组为基础解系的齐次线性方程组。
\end{itemize}

\subsection{题目5}
5．（15 分）设 \(\mathbb{C}[x]_{n^{2}}\) 是复数域上所有次数小于 \(n^{2}\) 的多项式以及零多项式构成的集合，对于通常的多项式加法和数乘，判断 \(\mathbb{C}[x]_{n^{2}}\) 能否构成实数域 \(\mathbb{R}\) 上的线性空间并简要证明你的结论。如果 \(\mathbb{C}[x]_{n^{2}}\) 是\(\mathbb{R}\) 上的线性空间，求出它的维数和一组基。

\subsection{题目6}
6．（15 分）设线性空间 \(V\) 上的线性变换 \(\sigma\) 在 \(V\) 的一组基 \(\varepsilon_{1}, \varepsilon_{2}, \varepsilon_{3}, \varepsilon_{4}\) 下的矩阵是

\[
A=\left(\begin{array}{cccc}
1 & 0 & 0 & 0 \\
-1 & -1 & -1 & 0 \\
1 & 1 & 1 & 0 \\
4 & 4 & 4 & -1
\end{array}\right)
\]

求 \(V\) 的一组基，使得 \(\sigma\) 在这组基下的矩阵是若尔当形矩阵。

\subsection{题目7}
7．（15 分）设矩阵 \(A\) 是实数域 \(\mathbb{R}\) 上的 \(n\) 阶可逆矩阵，证明：矩阵 \(A\) 能唯一地分解为 \(A=Q U\) ，其中\(Q\) 是正交矩阵，\(U=\left(u_{i j}\right)_{n \times n}\) 是上三角矩阵且对角元 \(u_{i i}(1 \leq i \leq n)\) 均为大于 0 的实数．

\subsection{题目8}
8．（20 分）设 \(a_{1}, a_{2}, \cdots, a_{m}\) 是数域 \(P\) 上的 \(m\) 个不同的数，\(b_{1}, b_{2}, \cdots, b_{m}\) 是数域 \(P\) 上的任意 \(m\) 个数，对于任一多项式 \(p(x)\) ，我们记 \(\partial(p(x))\) 为其次数。证明：存在唯一的多项式 \(g(x) \in P[x]\) ，使得\(\partial(g(x))<m\) ，且对任意的 \(1 \leq i \leq m\) ，我们有 \(g(x)=q_{i}(x)\left(x-a_{i}\right)+b_{i}\) ，其中 \(q_{i}(x) \in P[x]\) ．

\subsection{题目9}
9．（20 分）设实数域上的 \(n\) 阶方阵 \(A\) 满足 \(A^{2}=O\) ，记

\[
B=A A^{\mathrm{T}}+A^{\mathrm{T}} A, C=A+A^{\mathrm{T}} .
\]

令 \(\operatorname{Ker}(B), \operatorname{Ker}(C), \operatorname{Im}(A)\) 和 \(\operatorname{Im}\left(A^{\mathrm{T}}\right)\) 为如下定义的实向量空间 \(\mathbb{R}^{n}\) 的子空间：

\[
\begin{aligned}
& \operatorname{Ker}(B)=\left\{X \in \mathbb{R}^{n} \mid B X=0\right\}, \operatorname{Ker}(C)=\left\{X \in \mathbb{R}^{n} \mid C X=0\right\} \\
& \quad \operatorname{Im}(A)=\left\{A X \mid X \in \mathbb{R}^{n}\right\}, \operatorname{Im}\left(A^{\mathrm{T}}\right)=\left\{A^{\mathrm{T}} X \mid X \in \mathbb{R}^{n}\right\}
\end{aligned}
\]

证明：

\begin{itemize}
  \item[（1）] \(\operatorname{Ker}(B)=\operatorname{Ker}(C)\) ．
  \item[（2）] \(\mathbb{R}^{n}=\operatorname{Im}(A) \oplus \operatorname{Im}\left(A^{\mathrm{T}}\right) \oplus \operatorname{Ker}(B)\) ．
\end{itemize}
\end{document}
